% Kamp III.1: source-specific normal form over Lin.

\paragraph{The first-order normal form over \(\mathsf{Lin}\).}
The following construction gives the preliminary result of Kamp's
Theorem~III.1 and its sentence corollary.  The order remains an arbitrary
nonempty strict linear order, with arbitrary interpretations of the unary
predicates.  All intervals below are open; they can be empty.  Neither
density, absence of endpoints, nor order completeness is assumed
\cite[ch.~III, defs.~4--9 and thm.~1, pp.~43--56]{Kamp1968}.

\begin{hbdefinition}{def:p09-kamp-iii1-prefix}{Kamp's temporal prefixes and their representatives}
Work in the equality-free language \(L\) already specified.  In this
construction the source's displayed equality is represented by the
abbreviation
\[
 x\doteq y:=\neg(x<y)\land\neg(y<x).
\]
It expresses coincidence of the assigned points over \(\mathsf{Lin}\).
It is not an additional predicate symbol.  Also put
\[
 \top_x:=\neg(x<x),\qquad \bot_x:=x<x.
\]
These are respectively true and false at every point of every order in
\(\mathsf{Lin}\), and both retain the displayed free variable.

A \emph{temporal prefix} is a finite conjunction of comparisons
\(x<y\) or \(x\doteq y\).  It is \emph{complete} on its variables if it
decides every such comparison over \(\mathsf{Lin}\), and
\emph{consistent} if some linear-order assignment satisfies it.
Order variable names by their indices in \(t_0,t_1,\ldots\).
For a finite nonempty set \(X\) of variables, a complete consistent prefix
specifies an ordered partition
\[
 (B_1,\ldots,B_s)\quad\text{of }X.
\]
Variables in one block denote the same point, and every point of \(B_k\)
precedes every point of \(B_{k+1}\).  Let \(y_k\) be the first variable
name in \(B_k\).  These are Kamp's \emph{essential variables}; their
\emph{order in the prefix} is \(y_1,\ldots,y_s\), the order of their
assigned points, which need not agree with their variable-name order.

The \emph{standard prefix} for this partition contains the chain
\(y_1<\cdots<y_s\), followed by
\(y_k\doteq v\) for each \(v\in B_k\setminus\{y_k\}\).
The chain is omitted when \(s=1\).
Write the equalities block by block, and within each block in variable-name
order.  When \(X=\{y_1\}\), use \(y_1\doteq y_1\).
This writes only adjacent strict comparisons, puts strict comparisons before
equalities, and uses representatives on the left of equalities; a redundant
variable occurs only on the right of its one equality.  It is the source's
standard-prefix convention with an explicit ordering of same-block
equalities and of parentheses.  The number \(s\) is the source's
\(l(\pi)\); \(y_1\) and \(y_s\) are its first and last essential
variables.  The prefix does not assert the existence or absence of other
points between them.

Before a construction, rename bound variables so that no variable is both
free and bound or bound by two quantifiers.  Kamp additionally orders free
variables before bound variables and outer bound variables before inner
ones when speaking of precedence in a formula.  In a temporal prefix all
its variables are free, so representative selection uses the variable-name
order just specified.  Every substitution below concerns free occurrences
and is capture-avoiding.
\end{hbdefinition}

\begin{hbdefinition}{def:p09-kamp-iii1-components}{The \(\operatorname{At}\), \(\operatorname{Bef}\), \(\operatorname{Aft}\), and \(\operatorname{Bet}\) components}
Kamp's literal class \(\operatorname{At}(x)\) consists of finite nonempty
conjunctions of unary literals \(P_i(x)\) and \(\neg P_i(x)\), together
with the particular formula \(P_0(x)\lor\neg P_0(x)\).
A literal is an atomic unary formula or its negation.  Write
\[
 T^K_x:=P_0(x)\lor\neg P_0(x),\qquad
 F^K_x:=P_0(x)\land\neg P_0(x).
\]
Both belong to the source's \(\operatorname{At}(x)\).  The order formulas
\(\top_x,\bot_x\) of Definition~\ref{def:p09-kamp-iii1-prefix} do not
belong to that unextended syntactic class.

Initially
\[
 \operatorname{Bef}_0(x)=\operatorname{Aft}_0(x)
 =\operatorname{Bet}_0(x,y)=\varnothing.
\]
For each row of the following display, retain all components of stage \(n\)
and add the two displayed forms, where \(m\geq1\), each \(\theta_i\)
belongs to the indicated set, and \(z\) is a fresh variable distinct from
the endpoints:
\[
 \exists z\bigl(\beta\land\bigwedge_{i=1}^{m}\theta_i\bigr),
 \qquad
 \forall z\bigl(\beta\longrightarrow\bigvee_{i=1}^{m}\theta_i\bigr).
\]
\[
\begin{array}{c|c|l}
 \text{new class}&\beta&\text{permitted components }\theta_i\\ \hline
 \operatorname{Bef}_{n+1}(x)&z<x&
 \operatorname{At}(z)\cup\operatorname{Bef}_n(z)
                    \cup\operatorname{Bet}_n(z,x)\\
 \operatorname{Aft}_{n+1}(x)&x<z&
 \operatorname{At}(z)\cup\operatorname{Aft}_n(z)
                    \cup\operatorname{Bet}_n(x,z)\\
 \operatorname{Bet}_{n+1}(x,y)&x<z\land z<y&
 \operatorname{At}(z)\cup\operatorname{Bet}_n(x,z)
                    \cup\operatorname{Bet}_n(z,y)
\end{array}
\]
The unindexed classes are the unions over finite stages \(n\).
In the literal source grammar, an empty conjunction inside an existential
constructor is written as \(T^K_z\), and an empty disjunction inside a
universal constructor as \(F^K_z\).  These are formulas in the source's
fixed language \(L\), which contains \(P_0\).
The \(\operatorname{Bef}\) and \(\operatorname{Aft}\) components have only
their one endpoint free, and a \(\operatorname{Bet}\) component has only
its two displayed endpoints free.  The constructors describe nested
first-order conditions on regions; they are not yet temporal operators.

For a separate predicate-support-preserving rendition, define
\(\operatorname{At}^{*}(x)\) to contain the unary-literal conjunctions
and the two distinguished constants \(\top_x,\bot_x\), and repeat the
same recursion with every component class starred.  In that recursion use
\(\top_z\) and \(\bot_z\) for the two empty-list cases.  These
starred classes are an explicitly different syntactic presentation, not a
claim that an order atom is a source \(\operatorname{At}\) component.
The syntactic map \(\varepsilon\) replaces each distinguished starred
truth/falsehood slot by \(T^K_x/F^K_x\) at the same variable and commutes
with the regional constructors.  It maps the starred classes into the
unstarred source classes and preserves truth over \(\mathsf{Lin}\).
\end{hbdefinition}

\begin{hbdefinition}{def:p09-kamp-iii1-standard}{Standard formulas and standard sentences}
For a standard prefix \(\pi\) on \(X\), with essential variables
\(y_1,\ldots,y_s\), let
\[
\begin{split}
 \mathcal C_\pi={}&
 \bigcup_{k=1}^{s}\operatorname{At}(y_k)
 \ \cup\ \bigcup_{k=1}^{s-1}\operatorname{Bet}(y_k,y_{k+1})\\
 &{}\cup\operatorname{Bef}(y_1)\cup\operatorname{Aft}(y_s).
\end{split}
\]
A \emph{standard clause} on \(X\) is
\(\pi\land\bigwedge_j\theta_j\), with every
\(\theta_j\in\mathcal C_\pi\).  If it has no local condition, insert
\(T^K_{y_1}\).  A \emph{standard formula} on \(X\) is a finite
disjunction of standard clauses.  An always-false result can be written by
putting \(F^K_{y_1}\) in a clause for each standard prefix on \(X\).
Thus its named free variables are preserved as well.

A \emph{standard sentence} is \(\exists x\,N(x)\) or
\(\forall x\,N(x)\), where \(N\) is standard on \(\{x\}\).
These are Kamp's Definitions~8--9
\cite[ch.~III, pp.~48, 56]{Kamp1968}.  A \emph{starred standard formula}
or \emph{starred standard sentence} uses the starred component classes and
\(\top_x,\bot_x\) in the same positions.  Its image under
\(\varepsilon\) is a standard formula or sentence in the literal source
grammar.
\end{hbdefinition}

\begin{hbtheorem}{thm:p09-kamp-iii1}{proved=annex:p9.kamp.iii1}{Kamp's linear-order normal form and sentence corollary}
Let \(X=\{x_1,\ldots,x_r\}\), with \(r>0\), be the free variables of
an \(L\)-formula \(\varphi\).  There is an effectively constructed
standard formula \(N_\varphi\) on \(X\) such that
\[
 \mathsf{Lin}\models\varphi\longleftrightarrow N_\varphi.
\]
Equivalence holds for every assignment of \(X\) and every interpretation
of the unary predicates of \(L\).  The source grammar permits its fixed
\(P_0\) in constant components even if the input did not contain it.
Every sentence of \(L\) is \(\mathsf{Lin}\)-equivalent to a Boolean
combination of standard sentences.

\emph{Additional handbook property.}  There is also an effectively
constructed starred standard formula \(N^*_{\varphi}\), with exactly the
same free variables and no additional unary predicate, equivalent to
\(\varphi\) over \(\mathsf{Lin}\).  Taking
\(N_{\varphi}=\varepsilon(N^*_{\varphi})\) gives the literal
source-standard result above.  Predicate-support preservation is a property
of the starred rendition; it is not an additional restriction silently
imposed on Kamp's syntactic classes.

The construction first puts the input into prenex form.  It normalizes the
quantifier-free matrix by complete order types, then removes the original
quantifiers from the outside of the already normalized inner formula, using
the induction in Annex~\ref{annex:p9.kamp.iii1}.  For an existential
quantifier it retains one of the equal-point, left-ray, right-ray, or
intervening-gap possibilities.  For a universal quantifier it groups clauses
by the surviving prefix and by all possible positions of the quantified
point, including positions for which the matrix has no disjunct.  All
required cases, including empty regions, are proved in the annex
\cite[ch.~III, thm.~1 and cor.~1, pp.~48--56]{Kamp1968}.
\end{hbtheorem}

\paragraph{Witness existence at an endpoint or in a gap.}
The starred standard form of \(\exists x\,(x<y\land P(y))\) is
\[
 (y\doteq y)\land P(y)\land
 \exists x\,(x<y\land\top_x).
\]
The last component belongs to \(\operatorname{Bef}^{*}_1(y)\); its
\(\varepsilon\)-image belongs to \(\operatorname{Bef}_1(y)\).
On a one-point order with \(P(y)\) true the formula is false; on
\(\{0,1\}\) with \(y=1\) and \(P(1)\) true it is true.
The mirror \(\exists x\,(y<x\land P(y))\) uses
\(\operatorname{Aft}^{*}_1(y)\) and is false when \(y\) is the last point.
For a gap the component
\[
 (y<z)\land\exists x\,(y<x\land x<z\land\top_x)
\]
is false at consecutive integers \(y=0,z=1\), and true for the same
endpoints in the rational order.  A strict comparison between named
endpoints does not itself supply an intervening witness.

The universal counterpart has the form
\[
 \forall x\,(x<y\longrightarrow P(y))
 \quad\longleftrightarrow\quad
 P(y)\lor\forall x\,(x<y\longrightarrow\bot_x).
\]
The second disjunct records the empty left ray.  The analogous formulas with
\(y<x\) and \(y<x<z\) record a last point and an empty intervening gap.
These examples explain why the universal construction retains an emptiness
condition when its remaining disjuncts are independent of the quantified
variable.

\paragraph{Scope of this normal form.}
The output still contains nested first-order quantification in its
\(\operatorname{Bef}\), \(\operatorname{Aft}\), and
\(\operatorname{Bet}\) components.  Theorem~\ref{thm:p09-kamp-iii1}
establishes the preliminary \(\mathsf{Lin}\) normal form and the sentence
corollary.  Expressing the one-variable components in the image of \(D\)
is the further task of Theorem~III.2 and its completeness-dependent lemmas.
